% ====================================================================
% Introduction Section (Revised)
% ====================================================================
\section{Introduction and Motivation}\label{sec:introduction}

% --- Subsection 1.1: Combined Narrative ---
\subsection{The Case for Adaptive IoT Security}\label{subsec:case_for_adaptive_security}
The Internet of Things (IoT) represents a monumental shift in the digital landscape, weaving a network of interconnected devices into the fabric of modern society. This paradigm extends beyond consumer gadgets to critical infrastructure in healthcare, transportation, and industry \cite{Chen2021}. With the number of Internet of Things (IoT) devices worldwide being forecast to more than double from 19.8 billion in 2025 to more than 40.6 billion by 2034 \cite{Statista2025}, this exponential growth has created a vast and multifaceted attack surface. The financial ramifications are staggering, as global damages from cybercrime are estimated to reach USD 12.2 trillion annually by 2031 \cite{Morgan2025}. Malicious actors increasingly exploit insecure devices to steal data, disrupt services, or form large-scale botnets for Distributed Denial of Service (DDoS) attacks \cite{Neto2023, Tawalbeh2020}.

Securing this landscape presents unique and formidable challenges that render traditional IT security paradigms insufficient. First, most IoT devices are severely \textbf{resource-constrained}, possessing minimal computational power, memory, and energy, which precludes the use of intensive security software like traditional firewalls or antivirus agents \cite{Chen2021}. Second, IoT ecosystems are defined by extreme \textbf{heterogeneity}, comprising countless devices using disparate protocols and operating systems, which makes uniform security policy enforcement nearly impossible \cite{Mavroeidis2023}. This complexity is amplified by the large-scale, dynamic nature of these networks and the physical vulnerability of devices often deployed in unsecured locations \cite{AlGaradi2020}.

Consequently, conventional security tools are fundamentally ill-equipped for this environment. \textbf{Signature-based Intrusion Detection Systems (IDS)}, the cornerstone of many legacy systems, relies on static databases of known attack patterns. This approach is purely reactive and fails against novel or \textbf{zero-day attacks}—threats for which no signature yet exists \cite{Khraisat2019, Yang2024}. While \textbf{anomaly-based IDS} attempts to overcome this by baselining "normal" behavior, the dynamic and heterogeneous nature of IoT traffic makes establishing a reliable baseline exceptionally difficult, leading to high \textbf{false positive rates} that can disrupt critical functions \cite{Ibitoye2021, Tharewal2022}.

The core limitation of these paradigms is their lack of adaptability. They are built on static models and cannot autonomously evolve. This creates a critical need for a new security paradigm: one that is forward-looking, data-driven, and capable of autonomous decision-making. \textbf{Deep Reinforcement Learning (DRL)}, a field of machine learning focused on goal-oriented learning from interaction, has emerged as a powerful framework for such a system \cite{Chen2021}. A DRL agent learns an optimal decision-making strategy, or \textbf{policy}, through trial-and-error interaction with its environment. Unlike static systems, a DRL agent can generalize from its experience to respond to novel threats and can continuously adapt its policy as new adversarial tactics emerge, making it ideally suited for the non-stationary nature of IoT security \cite{Sutton2018, Uprety2021}. By framing cyber-defense as a sequential decision-making problem, DRL provides a robust mathematical foundation for creating autonomous security systems that can learn to outmaneuver attackers \cite{Gueriani2023}.

% --- Subsection 1.2: Renamed and Focused ---
\subsection{Proposal}\label{subsec:proposal_and_contributions}
The preceding analysis culminates in a clear research problem: a lack of security frameworks that are simultaneously \textbf{proactive}, \textbf{adaptive}, and \textbf{data-driven}—capable of anticipating future threats and learning to execute optimal defensive actions autonomously. Traditional reactive systems are insufficient for the dynamic IoT threat landscape.

This thesis directly addresses this gap by designing, implementing, and evaluating an adaptive defense system for IoT networks that leverages proactive attack prediction and Deep Reinforcement Learning. The primary goal is to validate a closed-loop framework that intelligently shifts from a reactive to a proactive defense posture. To this end, this work provides the following contributions:

\begin{enumerate}
    \item \textbf{A Novel Proactive-Adaptive Defense Framework:} We design and implement a two-stage architecture that uniquely combines a predictive LSTM-based model for threat anticipation with a DRL-based agent for autonomous action selection. This architecture establishes a synergistic feedback loop wherein threat prediction directly informs defensive action.
    
    \item \textbf{Grounding in a Realistic, Modern Dataset:} The system is developed and validated using the \textbf{CICIoT2023 dataset}. By leveraging this large-scale and recent dataset, we ensure our models are trained on data reflecting contemporary, real-world IoT attacks.
    
    \item \textbf{Comprehensive Benchmarking of DRL Algorithms:} We implement and conduct a rigorous comparative analysis of three state-of-the-art DRL agents: the value-based \textbf{Deep Q-Network (DQN)}, and the policy-gradient \textbf{Proximal Policy Optimization (PPO)} and \textbf{Advantage Actor-Critic (A2C)}. This benchmark provides crucial insights into the performance trade-offs of different RL paradigms in a cybersecurity context.
    
    \item \textbf{Development of a High-Fidelity Simulation Environment:} We create a custom, lightweight, and reproducible \textbf{Gymnasium environment} for training the DRL agents. A key innovation is the direct integration of the LSTM model's predictions into the agent's observation space, realistically simulating a proactive defense scenario.
\end{enumerate}

The remainder of this thesis is structured as follows. Section \ref{sec:background} provides a detailed review of background concepts and related work in the fields of IoT security and DRL-based defense. Section \ref{sec:methodology} details the proposed framework architecture.